\begin{abstract}
This report presents P1 of a virtual data integration project covering accommodation, restaurants and cultural heritage in Lisbon. Gonçalo Serafim develops guest reviews in Elasticsearch and accommodation listings, owners and price quotes in MongoDB; Tiago Vailant develops restaurants and simulated activity in PostgreSQL and the cultural catalogue in BaseX. The report describes the requirements, conceptual, logical and physical models, provenance, preparation and quality assessment of the four sources. The experiment reproducibly introduces missing values, invalid values and conflicting versions. Owners use synthetic Portuguese Citizen Card or British passport identifiers; activity and document validity are assessed separately. Twelve queries combine both students' subtopics, with explicit conditions, expected results and difficulty criteria. Measured local-source results support preparation of the P2 mediator.

\medskip\noindent\textbf{Keywords:} Virtual data integration, data modelling, heterogeneity, data quality, Lisbon.
\end{abstract}
